\chapter{Market-Risk Measures}\label{ch:rc:market-risk-measures}

In October 1994 J.P. Morgan published RiskMetrics: a method for turning a trading firm's
positions into one number, the loss it should not exceed on more than one day in twenty (or
a hundred), together with the daily volatilities and correlations needed to compute it, free
of charge. Two years later the Basel Committee let banks set their market-risk capital from
their own version of the number, provided they checked it every day against what actually
happened. \omterm{def:rc:market-risk-measures:var}{Value at risk} became the language of risk limits, capital and disclosure, and its
weaknesses became as well known: it says nothing about the size of the losses beyond it, it
can reward concentration, and its estimates lag the markets it measures. This chapter opens
Part~IV by computing the number three ways on a real book and real data, replacing it with
\omterm{def:rc:market-risk-measures:es}{expected shortfall}, backtesting both, and allocating them to their sources.

\section{Loss distributions and value at risk}

\begin{definition}[Risk factor]\label{def:rc:market-risk-measures:factor}
A \emph{risk factor}\index{risk factor} is a market variable whose changes drive the value
of positions: a yield at a tenor, an exchange rate, an implied volatility, a credit spread.
A risk model describes the joint distribution of risk-factor changes over a horizon and
revalues the positions under them.
\end{definition}

\begin{definition}[Value at risk]\label{def:rc:market-risk-measures:var}
The \emph{value at risk}\index{value at risk} of a portfolio at confidence level $\alpha$ over a
horizon $h$ is the $\alpha$-quantile of its loss $L = -\Delta V$ over $h$:
$\mathrm{VaR}_\alpha = \inf\{\ell:\P(L>\ell)\le1-\alpha\}$. At 99\% over one day it is a loss exceeded on
one day in a hundred.
\end{definition}

\begin{example}[The book]\label{ex:rc:market-risk-measures:book}
A small USD book on 23 September 2026: long USD~200 million of a ten-year Treasury par bond,
short USD~300 million of a two-year, long EUR~100 million, short JPY~5 billion, and short a
three-month at-the-money EURUSD straddle on EUR~400 million (implied volatility 8\%, held
fixed). Its four \omterm{def:rc:market-risk-measures:factor}{risk factors} are the daily changes in the two- and ten-year par yields (4.85\%
and 5.11\% that day) and the log returns of EURUSD (1.1411) and USDJPY (157.92). Their EWMA
daily volatilities are 6.4 and 5.7 basis points, 0.28\% and 0.65\%.
\end{example}

\section{Value at risk three ways}

\begin{definition}[Historical simulation]\label{def:rc:market-risk-measures:hs}
\emph{Historical simulation}\index{historical simulation} applies each of the last $n$
observed risk-factor changes to today's positions, revalues, and reads VaR from the empirical
distribution of the $n$ P\&Ls: with $k = \lceil(1-\alpha)n\rceil$, the $k$-th largest loss.
\end{definition}

\begin{definition}[Parametric value at risk]\label{def:rc:market-risk-measures:param}
\emph{Parametric value at risk}\index{parametric value at risk} (delta-normal) linearises the
P\&L in the factor changes, $\Delta V\approx\delta^\top x$, assumes $x$ normal with covariance $\Sigma$, and
gives $\mathrm{VaR}_\alpha = z_\alpha\sqrt{\delta^\top\Sigma\delta}$. With exponentially weighted covariances
($\Sigma_t = \lambda\Sigma_{t-1}+(1-\lambda)x_tx_t^\top$, $\lambda = 0.94$ for daily data in RiskMetrics) it
follows volatility quickly.
\end{definition}

\begin{definition}[Monte Carlo value at risk]\label{def:rc:market-risk-measures:mc}
\emph{Monte Carlo value at risk}\index{Monte Carlo value at risk} simulates factor changes
from a model (here normal with the EWMA covariance), revalues the portfolio fully on each
draw, and reads the quantile.
\end{definition}

\begin{definition}[Filtered historical simulation]\label{def:rc:market-risk-measures:fhs}
\emph{Filtered historical simulation}\index{filtered historical simulation} rescales each
past factor change by the ratio of today's volatility estimate to the estimate on its day,
$\tilde x_s = x_s\,\hat\sigma_{\text{today}}/\hat\sigma_s$, before applying it: it keeps the historical
shapes and dependence of the moves and today's volatility level.
\end{definition}

\begin{definition}[Delta--gamma approximation]\label{def:rc:market-risk-measures:dg}
The \emph{delta--gamma approximation}\index{delta--gamma approximation} replaces full
revaluation by the second-order expansion $\Delta V\approx\delta^\top x+\tfrac12x^\top\Gamma x$, with the
first and second derivatives of the portfolio in the \omterm{def:rc:market-risk-measures:factor}{risk factors}.
\end{definition}

\begin{example}[One book, several numbers]\label{ex:rc:market-risk-measures:methods}
Over the last 500 days, \omterm{def:rc:market-risk-measures:hs}{historical simulation} gives a one-day 99\% VaR of USD~2.10 million;
the EWMA delta-normal VaR is 1.58 million, the equal-weighted one 1.60 million, Monte Carlo
with full revaluation on the EWMA covariance 1.66 million, and \omterm{def:rc:market-risk-measures:fhs}{filtered historical simulation}
2.09 million (\cref{fig:rc:market-risk-measures:methods}). \omterm{def:rc:market-risk-measures:hs}{Historical simulation} with
delta--gamma revaluation gives 2.10 million, delta alone 1.91: the short straddle's gamma adds
losses in both tails. The normal methods are lower because the historical moves have fatter
tails than the normal.
\end{example}

\omcode{code/firm/varmodel/firm_varmodel.py}{23}{50}{Historical VaR and ES, EWMA covariance and volatility, and filtered historical scenarios.}\label{lst:rc:market-risk-measures:var}

\begin{definition}[Square-root-of-time rule]\label{def:rc:market-risk-measures:sqrt}
The \emph{square-root-of-time rule}\index{square-root-of-time rule} scales a one-day VaR to $h$
days by $\sqrt h$: exact for independent, identically distributed normal changes with zero
mean and fixed positions, an approximation otherwise.
\end{definition}

\section{Expected shortfall and coherence}

\begin{definition}[Expected shortfall]\label{def:rc:market-risk-measures:es}
The \emph{expected shortfall}\index{expected shortfall} at level $\alpha$ is the average loss in the
worst $1-\alpha$ of outcomes, $\mathrm{ES}_\alpha = \frac1{1-\alpha}\int_\alpha^1\mathrm{VaR}_u\,du$; historically, the average of
the $k$ largest losses.
\end{definition}

\begin{definition}[Coherent risk measure]\label{def:rc:market-risk-measures:coherent}
A risk measure $\rho$ is a \emph{coherent risk measure}\index{coherent risk measure} if it is
monotone, translation invariant ($\rho(L+c) = \rho(L)+c$), positively homogeneous ($\rho(\lambda L) = \lambda\rho(L)$
for $\lambda\ge0$) and subadditive ($\rho(L_1+L_2)\le\rho(L_1)+\rho(L_2)$). \omterm{def:rc:market-risk-measures:es}{Expected shortfall} is coherent;
VaR is not subadditive in general.
\end{definition}

\begin{example}[Two concentrated bonds]\label{ex:rc:market-risk-measures:subadd}
Each of two bonds defaults independently with probability 0.9\%, losing 100. Each alone has a
99\% VaR of zero; together, the probability of at least one default is 1.79\%, so the 99\% VaR
is 100: diversifying raised VaR. \omterm{def:rc:market-risk-measures:es}{Expected shortfall} at 99\% is 90 for each bond and 100.8 for
the pair, less than their sum of 180.
\end{example}

For normal P\&L, ES at 97.5\% equals VaR at 99\% almost exactly (2.338 against 2.326 standard
deviations), which is why the Basel Committee chose 97.5\% when it replaced VaR by ES in its
2019 market-risk standard. On the book, the EWMA ES at 97.5\% is USD~1.59 million, next to the
VaR of 1.58; historical ES is 2.34 million against a historical VaR of 2.10: fat tails show up
in ES first.

\begin{omfigure}
\begin{tikzpicture}
\begin{axis}[width=11.5cm,height=5.4cm,ybar,bar width=9pt,
  symbolic x coords={HS,EWMA,equal-wt,MC,FHS,HS delta-gamma},xtick=data,
  ylabel={USD million},
  tick label style={font=\small},label style={font=\small},
  ymin=0,ymax=2.8,grid=major,grid style={gray!25},
  legend columns=2,legend style={font=\footnotesize,at={(0.5,-0.2)},anchor=north,draw=none,fill=none,/tikz/every even column/.append style={column sep=8pt}},
  table/col sep=comma]
\addplot[fill=omDef!60,draw=omDef] table[x=method,y=var99]{figdata/rates-credit-risk/21-market-risk-measures/methods.csv};
\addplot[fill=omThm!60,draw=omThm] table[x=method,y=es975]{figdata/rates-credit-risk/21-market-risk-measures/methods.csv};
\legend{99\% VaR,97.5\% ES}
\end{axis}
\end{tikzpicture}

\omcaption{One-day 99\% VaR and 97.5\% ES of the book on 23 September 2026 by method. Normal
methods agree with each other and give lower numbers; historical and \omterm{def:rc:market-risk-measures:fhs}{filtered historical
simulation} see the fat tails, and their ES exceeds their VaR. Data: US Treasury par yields,
ECB reference rates; the chapter's tutorial.}\label{fig:rc:market-risk-measures:methods}
\end{omfigure}

\begin{omfigure}
\begin{tikzpicture}
\begin{axis}[width=11.5cm,height=5.2cm,ybar,bar width=4pt,
  xlabel={one-day P\&L under historical scenarios (USD million)},ylabel={days},
  tick label style={font=\small},label style={font=\small},
  xmin=-5.5,xmax=2.5,ymin=0,ymax=75,grid=major,grid style={gray!25},
  table/col sep=comma]
\addplot[fill=omDef!40,draw=omDef] table[x=pnl,y=count]{figdata/rates-credit-risk/21-market-risk-measures/hist.csv};
\draw[omThm,very thick] (axis cs:-2.097,0) -- (axis cs:-2.097,40) node[above,font=\footnotesize] {VaR};
\draw[omThm,very thick,dashed] (axis cs:-2.335,0) -- (axis cs:-2.335,28) node[above left,font=\footnotesize] {ES};
\end{axis}
\end{tikzpicture}

\omcaption{Distribution of the book's P\&L over the last 500 historical days applied to
today's positions, with the 99\% VaR and the 97.5\% ES. Data: US Treasury, ECB; the
chapter's tutorial.}\label{fig:rc:market-risk-measures:hist}
\end{omfigure}

\section{Backtesting the measure}

\begin{definition}[VaR exception]\label{def:rc:market-risk-measures:exception}
A \emph{VaR exception}\index{VaR exception} is a day on which the realised (or hypothetical)
loss exceeds the VaR computed the day before. A correct 99\% model has exceptions on 1\% of days,
independently.
\end{definition}

\begin{definition}[Kupiec test, Christoffersen test]\label{def:rc:market-risk-measures:tests}
The \emph{Kupiec test}\index{Kupiec test} compares the number of exceptions $x$ in $n$ days with
the expected $pn$ by the likelihood ratio $-2\ln[(1-p)^{n-x}p^x/((1-\hat p)^{n-x}\hat p^x)]$, $\hat p = x/n$,
chi-square with one degree of freedom. The \emph{Christoffersen test}\index{Christoffersen test}
checks independence: it compares the probability of an exception after an exception with the
probability after a quiet day.
\end{definition}

\begin{definition}[Traffic-light test]\label{def:rc:market-risk-measures:traffic}
The \emph{traffic-light test}\index{traffic-light test} classifies a 99\% VaR model by its
exceptions over 250 days: green for 0 to 4, yellow for 5 to 9 with increases of 0.40, 0.50,
0.65, 0.75 and 0.85 in the capital multiplication factor of 3, red for 10 or more (an increase
of 1). Under the 1996 rules, capital is the higher of the previous day's ten-day VaR and the factor times
its average over the preceding sixty business days.
\end{definition}

\omcode{code/firm/varmodel/firm_varmodel.py}{77}{108}{Kupiec's and Christoffersen's likelihood-ratio tests and the Basel traffic light.}\label{lst:rc:market-risk-measures:backtest}

\begin{example}[Four years of backtesting]\label{ex:rc:market-risk-measures:backtest}
Recomputing each model every day from 2 September 2022 to 23 September 2026, on the same book
revalued at each day's levels, against the next day's P\&L: \omterm{def:rc:market-risk-measures:hs}{historical simulation} has 10
exceptions in 1\,000 days (Kupiec p-value 1.00) and none in the last 250, green; the EWMA
delta-normal model has 28 (Kupiec likelihood ratio 22.0, p-value $3\times10^{-6}$) and 8 in the last
250, yellow with a 0.75 increase; \omterm{def:rc:market-risk-measures:fhs}{filtered historical simulation} has 8 and 3, green
(\cref{fig:rc:market-risk-measures:backtest}). None fails the independence test at 5\%.
\end{example}

\begin{omfigure}
\begin{tikzpicture}
\begin{axis}[width=11.5cm,height=5.6cm,
  xlabel={trading days from 2 September 2022},ylabel={USD million},
  tick label style={font=\small},label style={font=\small},
  xmin=0,xmax=1000,ymin=-4.5,ymax=4.5,grid=major,grid style={gray!25},
  legend columns=2,legend cell align=left,legend style={font=\footnotesize,at={(0.5,-0.25)},anchor=north,draw=none,fill=none,/tikz/every even column/.append style={column sep=8pt}},
  table/col sep=comma]
\addplot[omExo!60,only marks,mark=*,mark size=0.5pt] table[x=day,y=pnl]{figdata/rates-credit-risk/21-market-risk-measures/backtest.csv};
\addplot[omDef,thick] table[x=day,y=var_hs]{figdata/rates-credit-risk/21-market-risk-measures/backtest.csv};
\addplot[omMeth,thick] table[x=day,y=var_ewma]{figdata/rates-credit-risk/21-market-risk-measures/backtest.csv};
\addplot[omThm,only marks,mark=x,mark size=2.5pt,thick] table[x=day,y=pnl]{figdata/rates-credit-risk/21-market-risk-measures/exceptions.csv};
\legend{daily P\&L,$-$VaR historical,$-$VaR EWMA,EWMA exceptions}
\end{axis}
\end{tikzpicture}

\omcaption{Four years of the book's daily P\&L against the negative of the previous day's
99\% VaR from \omterm{def:rc:market-risk-measures:hs}{historical simulation} and from the EWMA delta-normal model; crosses mark the
EWMA model's exceptions. The EWMA line reacts to volatility quickly but runs too close to zero
for fat-tailed days. Data: US Treasury, ECB; the chapter's tutorial.}\label{fig:rc:market-risk-measures:backtest}
\end{omfigure}

\begin{dated}{2026-09}{VaR and ES in the rules}\label{dat:rc:market-risk-measures:basel}
RiskMetrics was launched by J.P. Morgan in October 1994 (fourth edition of its technical
document, with Reuters, December 1996), with decay factors of 0.94 for daily and 0.97 for
monthly horizons. The Basel Committee's 1996 backtesting framework set the traffic light for
99\% VaR over 250 days. Its 2019 market-risk standard replaced VaR by \omterm{def:rc:market-risk-measures:es}{expected shortfall} at
97.5\% on a ten-day base horizon, with desk-level backtests at 97.5\% and 99\%: more than 12
exceptions at 99\% or 30 at 97.5\% in twelve months send a desk to the standardised approach.
\end{dated}

\section{Risk contributions}

Parametric VaR is homogeneous of degree one in the positions, so \omterm{def:rc:the-valuation-adjustment-desk:euler}{Euler allocation}
(chapter~20) splits it into contributions that add up: $c_i = z_\alpha\delta_i(\Sigma\delta)_i/\sigma$.

\begin{example}[Where the risk sits]\label{ex:rc:market-risk-measures:euler}
Of the EWMA VaR of USD~1.58 million, the ten-year bond contributes 114\%, EURUSD 23\%, the
two-year short $-32\%$ and the yen short $-5\%$ (\cref{fig:rc:market-risk-measures:euler}): the
two-year hedges part of the ten-year's rate risk, and the negative contributions mark hedges.
\end{example}

\begin{omfigure}
\begin{tikzpicture}
\begin{axis}[width=11.5cm,height=5cm,ybar,bar width=20pt,
  symbolic x coords={2y yield,10y yield,EURUSD,USDJPY},xtick=data,
  ylabel={contribution (USD million)},
  tick label style={font=\small},label style={font=\small},
  ymin=-0.8,ymax=2.2,grid=major,grid style={gray!25},
  nodes near coords,every node near coord/.append style={font=\footnotesize,/pgf/number format/fixed,/pgf/number format/precision=2},
  table/col sep=comma]
\addplot[fill=omDef!60,draw=omDef] table[x=factor,y=contrib]{figdata/rates-credit-risk/21-market-risk-measures/euler.csv};
\end{axis}
\end{tikzpicture}

\omcaption{Euler contributions of each \omterm{def:rc:market-risk-measures:factor}{risk factor} to the book's EWMA 99\% VaR. They add up
to the VaR; negative contributions come from positions that hedge the rest. Data: the chapter's
tutorial.}\label{fig:rc:market-risk-measures:euler}
\end{omfigure}

\section{Tutorial: VaR on public data}\label{tut:rc:market-risk-measures:tutorial}

\begin{tutorial}
\textbf{Goal.} Compute historical, parametric, Monte Carlo and filtered VaR and ES of the book
on Treasury and ECB data, backtest them over four years, and allocate the VaR. \textbf{End state:}
the numbers of
\cref{ex:rc:market-risk-measures:methods,ex:rc:market-risk-measures:backtest,ex:rc:market-risk-measures:euler}
and the five charts.
\begin{tutsteps}
\item \textbf{Data}: \texttt{data()} joins the par yields and the ECB rates; \texttt{changes()}
gives the factor moves.
\item \textbf{Measures}: \texttt{measures()} runs every method on the last 500 days.
\item \textbf{Backtest}: \texttt{backtest()} rolls the models over 1\,000 days.
\item \textbf{Charts}: \texttt{fig\_rc\_var.py}.
\end{tutsteps}
\textbf{What to change next.} Use $\lambda = 0.97$; drop the straddle and compare delta and full
revaluation; shorten the window to 250 days and watch the historical VaR react.
\end{tutorial}

\section{Build: the VaR model}\label{bld:rc:market-risk-measures:varmodel}

\begin{build}
\bfield{Purpose} The firm's market-risk measures and their tests, used by the stress,
capital and risk-engine chapters that follow.
\bfield{Interface} \texttt{hs\_var\_es(pnl, alpha)}; \texttt{ewma\_cov}, \texttt{ewma\_vol\_path},
\texttt{fhs\_scenarios}; \texttt{parametric\_var\_es}, \texttt{euler\_var}; \texttt{mc\_var\_es(revalue,
cov, alpha)}; \texttt{delta\_gamma}; \texttt{kupiec}, \texttt{christoffersen},
\texttt{traffic\_light}; \texttt{sqrt\_time}.
\bfield{Rules} Losses positive; historical VaR is the $\lceil(1-\alpha)n\rceil$-th largest loss and ES the
average of that many; EWMA weights normalised; traffic light for 250 days.
\bfield{Acceptance tests} \texttt{code/firm/varmodel/tests/}: VaR and ES of a known sample;
Monte Carlo against parametric for a linear book; Euler contributions add up; normal ES at
97.5\% near VaR at 99\%; Kupiec zero at the expected count; traffic-light boundaries;
Christoffersen detects clustering; filtered scenarios carry today's volatility.
\bfield{Stretch} Extreme-value tails; GARCH filtering; ES backtests; stressed calibration;
component ES.
\end{build}

\begin{omsources}
\item J.P.\ Morgan/Reuters, \emph{RiskMetrics --- Technical Document}, 4th ed., 1996.
\item P.\ Artzner, F.\ Delbaen, J.-M.\ Eber and D.\ Heath, ``Coherent measures of risk'',
\emph{Mathematical Finance} 9(3), 1999, 203--228.
\item P.\ Kupiec, ``Techniques for verifying the accuracy of risk measurement models'',
\emph{Journal of Derivatives} 3(2), 1995, 73--84.
\item P.\ Christoffersen, ``Evaluating interval forecasts'', \emph{International Economic Review}
39(4), 1998.
\item Basel Committee on Banking Supervision, backtesting framework (1996); \emph{Minimum capital
requirements for market risk} (2019).
\end{omsources}

\section{Exercises}

\begin{exercise}[$\star$]\label{exo:rc:market-risk-measures:1}
With 500 historical days, which loss is the 99\% VaR and which losses make up the 97.5\% ES?
\end{exercise}

\begin{exercise}[$\star$]\label{exo:rc:market-risk-measures:2}
Scale the EWMA one-day VaR to ten days with the \omterm{def:rc:market-risk-measures:sqrt}{square-root-of-time rule}. When is the rule
wrong?
\end{exercise}

\begin{exercise}[$\star$]\label{exo:rc:market-risk-measures:3}
A model has 7 exceptions in 250 days. Give the zone and the multiplication factor.
\end{exercise}

\begin{exercise}[$\star\star$]\label{exo:rc:market-risk-measures:4}
Verify the numbers of \cref{ex:rc:market-risk-measures:subadd}.
\end{exercise}

\begin{exercise}[$\star\star$]\label{exo:rc:market-risk-measures:5}
Why does the short straddle make delta-only VaR too low?
\end{exercise}

\begin{exercise}[$\star\star$]\label{exo:rc:market-risk-measures:6}
Why did the EWMA model fail the \omterm{def:rc:market-risk-measures:tests}{Kupiec test} while tracking volatility well?
\end{exercise}

\begin{exercise}[$\star\star\star$]\label{exo:rc:market-risk-measures:7}
\emph{Coding.} Compute the Kupiec statistic and p-value for 28 exceptions in 1\,000 days at 1\%.
\end{exercise}

\begin{exercise}[$\star\star\star$]\label{exo:rc:market-risk-measures:8}
\emph{Find the flaw.} ``Our historical VaR had no exception in the last 250 days, so it is too
conservative and we should switch to the lower EWMA number.''
\end{exercise}

\section{Problem: The Quarter-Past-Four Report}

\begin{problem}[{Weekend problem --- one number, and what it costs}]\label{pb:rc:market-risk-measures:1}
The book of \cref{ex:rc:market-risk-measures:book} is a desk whose capital is set by its EWMA
delta-normal model under the 1996 framework: three (plus the backtesting increase) times the
ten-day 99\% VaR, scaled from one day.

\medskip\noindent\textbf{Part I --- The number.}
\begin{enumerate}
\item Give the desk's one-day 99\% VaR and 97.5\% ES from the EWMA model and from \omterm{def:rc:market-risk-measures:hs}{historical
simulation}.
\item Which \omterm{def:rc:market-risk-measures:factor}{risk factor} dominates, and by how much?
\item Give the ten-day VaR.
\item Why is the ES barely above the VaR in the EWMA model?
\item Which number would you report to the board, and why?
\end{enumerate}

\medskip\noindent\textbf{Part II --- The backtest.}
\begin{enumerate}[resume]
\item Give the exceptions over the last 250 days and the zone.
\item Give the multiplication factor.
\item Give the capital with and without the increase.
\item How many exceptions did \omterm{def:rc:market-risk-measures:hs}{historical simulation} have in the same year?
\item Would \omterm{def:rc:market-risk-measures:fhs}{filtered historical simulation} change the capital?
\end{enumerate}

\medskip\noindent\textbf{Part III --- The model.}
\begin{enumerate}[resume]
\item Why does a normal model with fast volatility still have too many exceptions?
\item What would the 2019 standard require instead of this VaR?
\item How would the desk fare under its 12- and 30-exception rule?
\item What does the independence test add?
\item Which model would you propose, and what would it cost in capital?
\end{enumerate}

\medskip\noindent\textbf{Part IV --- Judgement.}
\begin{enumerate}[resume]
\item Why is the lowest VaR not the best model?
\item Can a desk game a VaR limit? How?
\item What should the quarter-past-four report contain besides VaR?
\item State the \emph{named result}: the desk's 99\% VaR and 97.5\% ES, its exceptions over 250
days and the capital multiplier they imply.
\item In one sentence: what does VaR measure and what does it not?
\end{enumerate}
\end{problem}

\section{Interview questions}

\begin{interviewq}[$\star$ \iqroles{risk, trader}]\label{iq:rc:market-risk-measures:1}
Define VaR and ES. Why did regulators move from one to the other?
\end{interviewq}

\begin{interviewq}[$\star\star$ \iqroles{risk}]\label{iq:rc:market-risk-measures:2}
Compare \omterm{def:rc:market-risk-measures:hs}{historical simulation}, parametric and Monte Carlo VaR.
\end{interviewq}

\begin{interviewq}[$\star\star$ \iqroles{researcher, risk}]\label{iq:rc:market-risk-measures:3}
Show that VaR is not subadditive with an example.
\end{interviewq}

\begin{interviewq}[$\star\star$ \iqroles{risk, bank}]\label{iq:rc:market-risk-measures:4}
How do you backtest a VaR model? What do the Kupiec and \omterm{def:rc:market-risk-measures:tests}{Christoffersen tests} check?
\end{interviewq}

\begin{interviewq}[$\star\star\star$ \iqroles{developer}]\label{iq:rc:market-risk-measures:5}
A historical VaR on 100\,000 trades with full revaluation takes six hours. How do you make it
fast?
\end{interviewq}

\begin{interviewq}[$\star\star\star$ \iqroles{researcher}]\label{iq:rc:market-risk-measures:6}
How would you backtest \omterm{def:rc:market-risk-measures:es}{expected shortfall}?
\end{interviewq}
